\documentclass[11pt]{article}

\usepackage[margin=1in]{geometry}
\usepackage{microtype}
\usepackage{enumitem}
\usepackage{xcolor}
\usepackage{hyperref}
\usepackage{listings}

\definecolor{accent}{HTML}{2563EB}
\definecolor{codebackground}{HTML}{F3F4F6}

\hypersetup{
  pdftitle={Spectra Project Documentation},
  pdfauthor={Xiang Yangcheng},
  colorlinks=true,
  linkcolor=accent,
  urlcolor=accent
}

\setlength{\parindent}{0pt}
\setlength{\parskip}{0.65em}
\setlist{itemsep=0.25em, topsep=0.4em}
\lstset{
  basicstyle=\ttfamily\small,
  backgroundcolor=\color{codebackground},
  breaklines=true,
  columns=fullflexible,
  frame=single
}

\title{
  \textbf{Spectra}\\[0.35em]
  \large A Graphics Research Workspace for Static and Live Scenes
}
\author{Xiang Yangcheng}
\date{\href{https://github.com/Xayah-Graphics/spectra}
  {github.com/Xayah-Graphics/spectra}}

\begin{document}
\maketitle

\begin{abstract}
Spectra is a C++23 graphics research workspace for inspecting
static PBRT scenes and live plugin-driven scenes. It maintains one
shared scene representation that can be consumed by interactive and
path-traced renderer backends.
\end{abstract}

\section{Design Goals}

\begin{itemize}
  \item Keep scene loading outside individual renderer backends.
  \item Share one active scene between visualization and rendering systems.
  \item Support static PBRT content and live plugin-provided scenes.
  \item Keep simulation and reconstruction projects independent from Spectra.
\end{itemize}

\section{System Architecture}

\subsection{Application Shell}

The application shell owns the window, Vulkan lifecycle, workspace,
scene selection, and renderer integration.

\subsection{Shared Scene Model}

The \texttt{spectra.scene} model stores the active cameras, geometry,
materials, volumes, lights, and debug attachments.

\subsection{Scene Inputs}

Static PBRT files and dynamic scene plugins are converted into the
same shared scene representation before rendering.

\section{Renderer Backends}

\subsection{Vulkan Rasterizer}

The rasterizer provides interactive previews, camera visualization,
meshes, point clouds, volumes, selection, and debug overlays.

\subsection{CUDA/OptiX Path Tracer}

The path tracer consumes the same active scene and produces
path-traced output using CUDA and NVIDIA OptiX.

\section{Scene Workflow}

\begin{enumerate}
  \item Open a PBRT file or a dynamic scene plugin.
  \item Resolve the source into the shared scene model.
  \item Select a renderer backend.
  \item Inspect or render the active scene.
\end{enumerate}

\section{Build and Run}

The current development toolchain uses C++23, CMake 4.4,
Vulkan 1.4, CUDA 13.2, NVIDIA OptiX 9.1, and Ninja.

\begin{lstlisting}
cmake -S . -B build -G Ninja \
  -DCMAKE_BUILD_TYPE=Release \
  -DSPECTRA_OPTIX_PATH=/path/to/OptiX-SDK

cmake --build build --parallel
./build/spectra
\end{lstlisting}

A scene can be opened directly:

\begin{lstlisting}
./build/spectra --scene /path/to/scene.pbrt
./build/spectra --scene /path/to/plugin.dll
\end{lstlisting}

\section{Project Links}

Source code:
\href{https://github.com/Xayah-Graphics/spectra}
{github.com/Xayah-Graphics/spectra}

\end{document}
